\documentclass[11pt,a4paper]{article}
\usepackage[utf8]{inputenc}
\usepackage[T1]{fontenc}
\usepackage{amsmath,amssymb,amsthm}
\usepackage{mathpazo}
\usepackage[margin=2.5cm]{geometry}
\usepackage{booktabs}
\usepackage{hyperref}
\usepackage{cleveref}
\usepackage{graphicx}
\usepackage{xcolor}

\definecolor{designed}{HTML}{5a8a6a}
\definecolor{destructive}{HTML}{c4756a}
\definecolor{present}{HTML}{8B4049}

\newtheorem{proposition}{Proposition}
\newtheorem{corollary}{Corollary}
\newtheorem{definition}{Definition}
\theoremstyle{remark}
\newtheorem{remark}{Remark}

\title{\textbf{A Formal Economic Model of Technological Sovereignty}\\[0.5em]
\large Analytical Appendix to ``Two Kinds of Sovereignty''}
\author{}
\date{}

\begin{document}
\maketitle

\begin{abstract}
This document presents a three-layer formal economic model that unifies the core mechanisms of the essay ``Two Kinds of Sovereignty'': the exploitation/exploration allocation tradeoff (March, 1991), path-dependent lock-in (Arthur, 1989), induced innovation under constraint (Hicks, 1932; Hayami \& Ruttan, 1971), and techno-economic paradigm windows (Perez, 2002). Layer~1 solves the social planner's optimal control problem analytically. Layer~2 validates the analytical results through agent-based evolutionary simulation. Layer~3 translates both into actionable policy instruments. The model is calibrated to European data and produces three principal outputs: optimal allocation ratios between present and future sovereignty, threshold conditions for when designed constraint dominates laissez-faire, and the shadow price of delay in contested technological domains.
\end{abstract}

\tableofcontents
\newpage

%% ============================================================
\section{Layer 1: Analytical Core}
%% ============================================================

\subsection{Setup}

A social planner allocates a unit resource flow $R$ between two activities at each moment $t$:

\begin{itemize}
    \item \textbf{Exploitation} $\alpha R$: secures and domesticates existing technology. Produces present sovereignty capital $K_p$.
    \item \textbf{Exploration} $(1-\alpha)R$: invests in new technological paradigms. Produces future sovereignty capital $K_f$.
\end{itemize}

\subsection{State Variables}

The system is characterised by four state variables:

\begin{center}
\begin{tabular}{lll}
\toprule
\textbf{Variable} & \textbf{Description} & \textbf{Range} \\
\midrule
$K_p(t)$ & Present sovereignty stock & $\geq 0$ \\
$K_f(t)$ & Future sovereignty stock & $\geq 0$ \\
$D(t)$ & Dependency level & $[0, 1]$ \\
$W(t)$ & Paradigm window openness & $[0, 1]$ \\
\bottomrule
\end{tabular}
\end{center}

\subsection{Dynamics}

\begin{definition}[Present Sovereignty Accumulation]
\begin{equation}
\frac{dK_p}{dt} = A \cdot (\alpha R)^{\kappa} - \delta_p K_p - \varphi(D) K_p
\label{eq:kp}
\end{equation}
where $A$ is total factor productivity for exploitation, $\kappa \in (0,1)$ imposes diminishing returns, $\delta_p$ is the base depreciation rate, and $\varphi(D) = \varphi_0 D^2$ is dependency-accelerated depreciation (convex in $D$).
\end{definition}

\begin{definition}[Future Sovereignty Accumulation]
\begin{equation}
\frac{dK_f}{dt} = W(t) \cdot \left[(1-\alpha)R\right]^{\beta} \cdot h(\sigma) - \delta_f K_f
\label{eq:kf}
\end{equation}
where $\beta \in (0,1)$ governs exploration returns, and $W(t)$ is the Perez window function. Crucially, exploration investment has zero marginal return once the window closes. For analytical tractability, Layer~1 treats the economy as targeting a single paradigm window at a time; the multi-domain portfolio is handled in Layer~3.
\end{definition}

\begin{definition}[Dependency Dynamics (Arthur Increasing Returns)]
\begin{equation}
\frac{dD}{dt} = \lambda(\bar{D} - D)D - \eta K_p D
\label{eq:D}
\end{equation}
This is a logistic process toward the lock-in ceiling $\bar{D}$, slowed by present sovereignty investment. The multiplicative $D$ term in $\eta K_p D$ ensures $D \in [0,1]$: when $D \to 0$, the reduction term vanishes, preventing negative dependency.
\end{definition}

\begin{definition}[Perez Window Dynamics]
\begin{equation}
W_j(t) = W_{0,j} \cdot \exp\left(-\gamma_j \cdot \max(0,\, t - t_{\text{open},j})\right) \quad \text{for } t \geq t_{\text{open},j}
\label{eq:window}
\end{equation}
Windows open at $t_{\text{open},j}$ with initial openness $W_{0,j}$ and close exponentially at rate $\gamma_j$.
\end{definition}

\subsection{The Constraint Instrument $\sigma$}

The essay's ``virtual export restriction'' enters the exploration production function as a non-monotonic modifier:

\begin{equation}
h(\sigma) = 1 + a\sigma \exp(-b\sigma), \quad a, b > 0
\label{eq:hsigma}
\end{equation}

\begin{proposition}[Hump Shape of $h(\sigma)$]
\label{prop:hump}
The function $h(\sigma)$ defined in \eqref{eq:hsigma} satisfies:
\begin{enumerate}
    \item $h(0) = 1$ (no constraint, baseline productivity)
    \item $h$ is maximised at $\sigma^* = 1/b$ with $h(\sigma^*) = 1 + a/(be)$
    \item $h(\sigma) \to 1$ as $\sigma \to \infty$
    \item $h(\sigma) > 1$ for all $\sigma > 0$
\end{enumerate}
\end{proposition}

\begin{proof}
$h'(\sigma) = a\exp(-b\sigma)(1 - b\sigma)$. Setting $h'(\sigma) = 0$ gives $\sigma^* = 1/b$. $h''(\sigma^*) = -ab\exp(-1) < 0$, confirming a maximum. For large $\sigma$, $a\sigma\exp(-b\sigma) \to 0$, so $h \to 1$. Since $a\sigma\exp(-b\sigma) > 0$ for all $\sigma > 0$, we have $h > 1$ everywhere.
\end{proof}

This hump shape formalises the essay's three selection pressure lanes:
\begin{itemize}
    \item \textbf{Absent pressure} ($\sigma = 0$): Baseline productivity, no induced innovation.
    \item \textbf{Designed pressure} ($\sigma \approx \sigma^*$): Constraint maximises exploration productivity.
    \item \textbf{Destructive pressure} ($\sigma \gg \sigma^*$): Excessive constraint degrades toward baseline.
\end{itemize}

\subsection{The Planner's Problem}

\begin{equation}
\max_{\alpha(t),\, \sigma(t)} \int_0^{\infty} e^{-\rho t} \cdot V(K_p, K_f, D)\, dt
\label{eq:planner}
\end{equation}

subject to \eqref{eq:kp}--\eqref{eq:D}, where the flow value is:

\begin{equation}
V(K_p, K_f, D) = \omega K_p + (1-\omega)K_f - \psi D^2
\label{eq:value}
\end{equation}

\begin{remark}[Rationale for the Linear-Quadratic Form]
Sovereignty benefits are linear because each unit of $K_p$ or $K_f$ provides roughly constant marginal strategic value. Dependency costs are quadratic because the damage is convex: going from 50\% to 60\% dependent is substantially worse than 20\% to 30\%, reflecting non-linear lock-in and narrowing exit options. This form ensures the optimal policy is smooth (interior solution) rather than bang-bang.
\end{remark}

\subsection{Analytical Outputs}

The model is solved numerically via forward simulation over a grid of $(\alpha, \sigma)$ values, with scalar optimisation refining the peak. Key outputs:

\begin{proposition}[Optimal Allocation]
Under the baseline calibration ($D_0 = 0.67$, $R = 0.022$, $\gamma = 0.1$), the optimal fixed allocation is $\alpha^* \approx 0.35$, with optimal constraint intensity $\sigma^* = 2.0$. This implies roughly 65\% of sovereignty-oriented resources should flow to exploration (future sovereignty) rather than exploitation (present sovereignty).
\end{proposition}

\begin{proposition}[Shadow Price of Delay]
The cost of delaying exploration investment by $\Delta t$ years is:
\begin{equation}
C(\Delta t) = V(\alpha^*, \sigma^*, W_0) - V(\alpha^*, \sigma^*, W_0 \cdot e^{-\gamma \Delta t})
\end{equation}
This cost is non-linear and accelerating: each additional year of delay costs more than the last, because windows close exponentially. Under baseline calibration, a 5-year delay costs approximately 3 times as much as a 1-year delay.
\end{proposition}

\subsection{Regime Comparison}

Three policy regimes produce strikingly different trajectories:

\begin{center}
\begin{tabular}{lrrrl}
\toprule
\textbf{Regime} & $\alpha$ & $\sigma$ & \textbf{$V$ (discounted)} & \textbf{Interpretation} \\
\midrule
Status Quo (Presentism) & 0.80 & 0.0 & $-2.55$ & Exploit-heavy, no constraint \\
\textcolor{designed}{Designed Pressure} & 0.35 & 2.0 & $3.93$ & Optimal allocation + constraint \\
\textcolor{destructive}{Overreaction} & 0.20 & 15.0 & $-5.09$ & Excessive constraint \\
\bottomrule
\end{tabular}
\end{center}

The designed pressure regime dominates both alternatives, yielding positive total value where both alternatives are negative. The status quo's negative value reflects the compounding cost of dependency under path-dependent lock-in.

%% ============================================================
\section{Layer 2: Evolutionary Simulation}
%% ============================================================

\subsection{Purpose}

The analytical core assumes a representative agent and smooth production functions. Layer~2 relaxes both, populating the economy with $N$ heterogeneous firms to capture emergent dynamics: why some constraint regimes produce compounding capability while others fragment the innovation system.

\subsection{Agent Structure}

Each of $N = 500$ firms is characterised by:

\begin{center}
\begin{tabular}{lll}
\toprule
\textbf{Attribute} & \textbf{Description} & \textbf{Initial Distribution} \\
\midrule
$c_i(t)$ & Capability stock & Log-normal, mean-normalised to 1 \\
$r_i(t)$ & Research orientation $[0,1]$ & Beta$(2, 5)$, exploitation-skewed \\
$x_i(t)$ & External dependency ratio & Beta$(5, 2)$, high-dependency-skewed \\
$s_i(t)$ & Market share & Uniform: $1/N$ \\
\bottomrule
\end{tabular}
\end{center}

\subsection{Innovation Process}

Each period ($\Delta t = 1$ year), firms draw from two innovation lotteries:

\paragraph{Exploitation.} Probability $p_0 \cdot c_i^{\nu}$, payoff $\varepsilon_{\text{exploit}} \cdot c_i \cdot x_i \cdot (1 - \sigma x_i)$. High probability, low variance. The constraint factor $(1 - \sigma x_i)$ means highly dependent firms face an immediate productivity shock under designed pressure.

\paragraph{Exploration.} Probability $q_0 \cdot c_i^{\nu} \cdot W(t)$, payoff drawn from $\text{Pareto}(\xi) \cdot \varepsilon_{\text{explore}} \cdot c_i \cdot (1 - x_i)$. Low probability, fat-tailed payoff. Exploration success also reduces the firm's dependency ratio.

\subsection{Selection Mechanism}

Market shares update each period via:
\begin{equation}
\tilde{s}_i(t+1) = s_i(t) \cdot \left(\frac{c_i(t)}{\bar{c}(t)}\right)^{\theta}, \quad s_i(t+1) = \frac{\tilde{s}_i(t+1)}{\sum_j \tilde{s}_j(t+1)}
\end{equation}

Firms with $s_i < s_{\min}$ exit and are replaced by new entrants drawn from the initial distributions.

\subsection{Emergent Regimes}

Under different $\sigma$ values, three qualitatively distinct system dynamics emerge, corresponding to the essay's three selection pressure lanes:

\begin{enumerate}
    \item \textbf{$\sigma = 0$ (Absent Pressure):} Firms drift toward high dependency. Aggregate capability grows slowly and undirected. Efficient but fragile.
    \item \textbf{$\sigma \approx \sigma^*$ (Designed Pressure):} High-capability firms redirect toward exploration. Aggregate capability dips initially then compounds as alternative pathways mature.
    \item \textbf{$\sigma \gg \sigma^*$ (Destructive Pressure):} Most firms lack the absorptive capacity to redirect. Mass substitution with inferior alternatives. Aggregate capability degrades.
\end{enumerate}

\subsection{Cross-Layer Validation}

Monte Carlo sweeps across $\sigma$ values confirm the hump-shaped relationship between constraint intensity and aggregate terminal capability, validating the analytical $h(\sigma)$ assumption. The simulation also reveals distributional dynamics invisible to Layer~1: the \emph{capability threshold} below which designed pressure destroys rather than redirects firm-level innovation.

\begin{proposition}[Capability Threshold]
There exists a minimum initial capability $c_{\min}(\sigma)$, increasing in $\sigma$, below which firms lose capability under designed pressure relative to the unconstrained baseline. This threshold determines the absorptive capacity precondition for successful strategic necessity programmes.
\end{proposition}

%% ============================================================
\section{Layer 3: Policy Design Interface}
%% ============================================================

\subsection{Observable State Variables}

The planner does not observe $K_f$ or $W(t)$ directly. The policy interface maps observable proxies to recommendations:

\begin{center}
\begin{tabular}{lll}
\toprule
\textbf{Observable} & \textbf{Proxy for} & \textbf{Empirical Anchor} \\
\midrule
$D_{\text{obs}}$ & Dependency $D(t)$ & Cloud market share: 65--70\% \\
$C_{\text{obs}}$ & Capability distribution & R\&D intensity: 2.2\% GDP \\
$W_{\text{obs}}$ & Window openness & Maturity curves, VC acceleration \\
\bottomrule
\end{tabular}
\end{center}

\subsection{Policy Instruments}

\paragraph{Instrument 1: Allocation Ratio $\alpha^*$.} How much sovereignty spending goes to exploitation vs.\ exploration. Given $(D_{\text{obs}}, C_{\text{obs}}, W_{\text{obs}})$, the model outputs $\alpha^*$ and its sensitivity to measurement error.

\paragraph{Instrument 2: Constraint Intensity $\sigma^*$.} The ``virtual export restriction'' strength, scaled by observed capability: $\sigma^* = \sigma^*_{\text{base}} \cdot C_{\text{obs}}$. Low-capability systems receive weaker designed pressure to avoid the destructive regime.

\paragraph{Instrument 3: Domain Selection.} Candidate paradigm windows scored on:
\begin{equation}
S_j = w_W \cdot W_{j,\text{remaining}} + w_C \cdot C_j + w_S \cdot V_j^{\text{strategic}}
\end{equation}

with weights $(w_W, w_C, w_S) = (0.4, 0.3, 0.3)$.

\subsection{Domain Assessments}

\begin{center}
\begin{tabular}{lcccc}
\toprule
\textbf{Domain} & $W_{\text{remaining}}$ & $C_j$ & $V_j^{\text{strategic}}$ & \textbf{Score} \\
\midrule
Fusion Energy & 0.8 & 0.5 & 0.9 & 0.74 \\
Scientific AI & 0.7 & 0.4 & 0.8 & 0.64 \\
Quantum & 0.6 & 0.6 & 0.7 & 0.63 \\
\bottomrule
\end{tabular}
\end{center}

Fusion scores highest due to a newly opened window and very high strategic value. Scientific AI and quantum are closely ranked, with different profiles: quantum has stronger existing European capability, while scientific AI has a wider remaining window.

%% ============================================================
\section{Calibration}
%% ============================================================

\begin{center}
\begin{tabular}{llll}
\toprule
\textbf{Parameter} & \textbf{Value} & \textbf{Source} & \textbf{Layer} \\
\midrule
$D_0$ & 0.67 & Synergy Research Group & 1, 2, 3 \\
$R$ & 0.022 & European Investment Bank & 1 \\
$\kappa$ & 0.6 & Cobb-Douglas convention & 1 \\
$\beta$ & 0.5 & Exploration returns & 1 \\
$\delta_p$ & 0.05 & Annual depreciation & 1 \\
$\delta_f$ & 0.03 & Knowledge decay & 1 \\
$\bar{D}$ & 0.95 & Lock-in ceiling & 1 \\
$\lambda$ & 0.3 & Lock-in speed & 1 \\
$a$ & 2.0 & Hump amplitude & 1 \\
$b$ & 0.5 & Hump decay $\Rightarrow \sigma^* = 2.0$ & 1 \\
$\omega$ & 0.5 & Equal weighting & 1 \\
$\rho$ & 0.03 & Social discount rate & 1 \\
$N$ & 500 & Firm count & 2 \\
$p_0$ & 0.7 & Exploitation success rate & 2 \\
$q_0$ & 0.15 & Exploration success rate & 2 \\
$\nu$ & 0.3 & Capability exponent & 2 \\
$\xi$ & 2.5 & Pareto tail parameter & 2 \\
$\theta$ & 0.5 & Selection intensity & 2 \\
\bottomrule
\end{tabular}
\end{center}

\subsection{Comparative Statics}

Sensitivity analysis reveals that the model's policy recommendations are most sensitive to:
\begin{enumerate}
    \item \textbf{Window closure rate $\gamma$}: Faster-closing windows dramatically increase the shadow price of delay and shift $\alpha^*$ toward exploration.
    \item \textbf{Lock-in speed $\lambda$}: Higher lock-in speed increases the urgency of present sovereignty investment, creating tension with the exploration imperative.
    \item \textbf{Exploration returns $\beta$}: Higher returns to exploration make the case for reallocation stronger, lowering $\alpha^*$.
\end{enumerate}

%% ============================================================
\section{Conclusion}
%% ============================================================

The model formalises the essay's central argument: Europe's sovereignty challenge is not simply a question of securing today's technology stack, but of allocating scarce resources between exploitation and exploration under closing paradigm windows and self-reinforcing dependency. The three-layer structure reveals that:

\begin{enumerate}
    \item The \textbf{optimal allocation} heavily favours exploration ($\alpha^* \approx 0.35$), suggesting current sovereignty spending is misallocated toward presentism.
    \item \textbf{Designed constraint} ($\sigma \approx \sigma^*$) can raise exploration productivity, but only when absorptive capacity exceeds a threshold --- programme design must take this seriously.
    \item The \textbf{cost of delay} is non-linear and accelerating. Each year of inaction narrows the remaining paradigm windows, compounding the eventual cost of building alternatives.
\end{enumerate}

Necessity is the mother of invention. But as this model shows, only \emph{designed} necessity --- calibrated to capability, timed to windows, and sustained by investment --- can serve as the grandmother of sovereignty.

\end{document}
